\section{Introduction : Le Défi des DSLs Industriels}

\subsection{Contexte : L'angle mort des Grands Modèles de Langage}

L'avènement des modèles de langage massivement pré-entraînés (LLM) a bouleversé les pratiques de développement logiciel. Sur des langages populaires comme Python ou Javascript, des outils comme GitHub Copilot agissent désormais comme des assistants quasi-omniscients. La raison de ce succès est simple : ces langages constituent une part significative du corpus d'entraînement public sur lequel ces modèles ont appris.

Cependant, la réalité industrielle est plus nuancée. De nombreuses entreprises, à l'instar de notre partenaire Lokad, fondent leur avantage compétitif sur des \textit{Domain Specific Languages} (DSL) propriétaires. Dans le cas de Lokad, il s'agit d'Envision \cite{LokadEnvision}, un langage proche du SQL dédié à l'optimisation de la Supply Chain. Ici réside une première partie de notre problématique : les LLM standards (GPT-4, Mistral, Claude) n'ont quasiment jamais "lu" de code Envision.

La principale difficulté du projet repose ainsi sur la complexité réelle des architectures Envision mises à disposition : un
compte typique est composé de plusieurs dizaines de scripts interconnectés, organisés en modules thématiques,
et dont la documentation est partielle. Ces scripts réalisent des tâches variées allant du traitement de données
à l’optimisation numérique.

\subsection{Objectifs}

Nous ne cherchons pas seulement à créer un moteur de recherche amélioré. Nous essayons de concevoir une \textbf{architecture agentique hybride}. Ce système doit être capable de décider de lui-même s'il doit effectuer une recherche sémantique (comprendre un concept) ou lexicale (trouver une variable exacte), et d'itérer sur sa propre réponse.

Nous détaillerons comment nous sommes passés d'une pipeline linéaire à l'utilisation de graphes d'états (via \textit{LangGraph}), comment nous avons surmonté les limites du découpage de code (\textit{chunking}), et comment nous avons mis en place un protocole d'évaluation automatisé (\textit{LLM-as-a-Judge}) pour mesurer nos progrès de manière objective.